\chapter{Installation and Setup Appendix}

\section{Choose an Installation Route}
uafR supports two student installation routes. General users can install the
current public source from GitHub. Classes and restricted-network settings can
use a versioned offline student bundle so every participant starts from the
same package archive and training files.

\subsection{Route A: Public GitHub Installation}
Use this route when the machine has reliable internet access and the program
wants the current public version.

\begin{rcode}
install.packages(c("remotes", "BiocManager"))

BiocManager::install(
  c("ChemmineR", "fmcsR"),
  ask = FALSE,
  update = FALSE
)

remotes::install_github(
  "castrattonDSU/uafR",
  dependencies = c("Depends", "Imports"),
  upgrade = "never",
  build_vignettes = FALSE
)

library(uafR)
packageVersion("uafR")
\end{rcode}

After installation, run the repository's offline training smoke test or the
student acceptance test supplied by the program. Installation from GitHub
requires internet access for both uafR and its dependencies.

\subsection{Route B: Versioned Offline Student Bundle}
Use this route when a class needs a fixed, auditable package archive or when
GitHub access is unavailable. The bundle installs uafR from a local source
archive and uses CRAN and Bioconductor only for dependencies that are not
already installed.

The student bundle includes:
\begin{itemize}
  \item \path{START_HERE.md}, the first checklist to open.
  \item \path{uafR_QUICK_REFERENCE.md}, the compact command reference.
  \item \path{CHECKSUMS.csv}, the generated file size and MD5 checksum table.
  \item \path{verify_bundle_integrity.R}, the checksum verification script.
  \item A local uafR source package archive in \path{packages/}.
  \item \path{preflight_check.R}, the machine readiness check.
  \item \path{install_uafR_from_bundle.R}, the student installer.
  \item \path{update_uafR_from_bundle.R}, the reinstall/update script for newer bundles.
  \item \path{verify_uafR_install.R}, the post-install verification script.
  \item \path{run_student_acceptance_test.R}, the offline readiness test.
  \item The current training manual PDF.
  \item Training scripts and data templates used during the internship.
\end{itemize}

The installer still needs internet access for missing CRAN and Bioconductor dependencies. The uafR package itself is installed from the local bundle archive.

\section{Install From the Student Bundle}
The installation workflow is:

\begin{enumerate}
  \item Install current R and RStudio.
  \item Unzip the whole bundle folder.
  \item Open RStudio.
  \item Open \path{START_HERE.md} from the unzipped bundle folder.
  \item Source \path{verify_bundle_integrity.R}.
  \item Source \path{preflight_check.R}.
  \item If preflight has no failed checks, source \path{install_uafR_from_bundle.R}.
  \item Source \path{run_student_acceptance_test.R}.
  \item Keep \path{uafR_QUICK_REFERENCE.md} nearby and open \path{training/uafR_training_manual.pdf}.
\end{enumerate}

The main setup commands are:

\begin{rcode}
source("verify_bundle_integrity.R")
source("preflight_check.R")
source("install_uafR_from_bundle.R")
source("run_student_acceptance_test.R")
\end{rcode}

The installer prints five visible steps: checking the bundle and R library, installing CRAN dependencies, installing Bioconductor dependencies, installing uafR from the local archive, and verifying the installation. The installer writes \path{uafR_install_log.txt} in the bundle folder.

\section{Confirm Bundle Integrity}
The bundle builder writes \path{CHECKSUMS.csv} with the file size and MD5 checksum for each file in the generated bundle. After a bundle is copied or unzipped, the integrity check confirms that the files still match the build record.

\begin{rcode}
source("verify_bundle_integrity.R")
\end{rcode}

It should end with:

\begin{verbatim}
Bundle integrity result: PASS.
\end{verbatim}

\section{Confirm the Installation Without Live Web Queries}
The verification script checks package loading, required dependencies, bundled datasets, and the saved standard categorate result. It does not query live web services.

\begin{rcode}
source("verify_uafR_install.R")
\end{rcode}

\section{Student Acceptance Test}
The acceptance test confirms that the installed package is ready for the training workflow. It runs package verification, loads the core data objects, and runs the offline workflow from \path{training/scripts/04_core_workflow.R}. It should end with:

\begin{verbatim}
Acceptance test result: PASS.
\end{verbatim}

\begin{rcode}
source("run_student_acceptance_test.R")
\end{rcode}

\section{Optional Live Smoke Test}
After the offline verification passes, run a two-compound live test when internet access is available:

\begin{rcode}
library(uafR)
data("library_data", package = "uafR")

quick_result <- categorate(
  compounds = c("aspirin", "caffeine"),
  chemical_library = library_data,
  input_format = "wide",
  detail = "research",
  cache = FALSE,
  throttle = 0.1,
  assay_detail_limit = 0
)

validateCategorateResult(quick_result)$Summary
quick_result$ChemicalTraitReport
\end{rcode}

The acceptance test can also include the live smoke test:

\begin{rcode}
Sys.setenv(UAFR_STUDENT_LIVE = "1")
source("run_student_acceptance_test.R")
\end{rcode}

\section{Species-First Plant Chemistry Extension}
The bundle includes \path{training/scripts/10_species_phytochemistry.R}. Its
default example uses clearly labeled simulated curated rows and does not make
live requests. It demonstrates the normalized plant result, validation,
evidence filtering, species matrix, candidate score, and export workflow.

\begin{rcode}
source("training/scripts/10_species_phytochemistry.R")
plant_example <- run_species_phytochemistry_smoke(
  output_dir = "results/plant_training_example"
)
\end{rcode}

The simulated rows teach package behavior; they are not evidence about a
student's samples and must not be reported as observed or published plant
chemistry. Live species discovery is an optional extension after the offline
workflow is understood.

\section{Troubleshooting}
\begin{tabularx}{\textwidth}{p{0.24\textwidth}p{0.32\textwidth}Y}
\toprule
\textbf{Problem} & \textbf{Likely cause} & \textbf{What to do}\\
\midrule
Bundle integrity check fails & One or more files changed, did not copy, or did not unzip correctly. & Delete the unzipped folder and unzip the original bundle again. If it still fails, get a fresh copy of the zip.\\
Installer cannot find uafR archive & The bundle was not fully unzipped or the \path{packages/} folder was moved. & Unzip the whole bundle again and keep the folder structure unchanged.\\
Preflight says the R library is not writable & R cannot install packages into the first library path. & Restart RStudio. If the problem remains, set a user library before installing.\\
Bioconductor dependency fails & Internet access, R version, or Bioconductor access is blocked. & Run \path{preflight_check.R}, update R if it is old, and try again on a network that can reach Bioconductor.\\
R asks about updating many packages & RStudio or Bioconductor is offering broad updates. & Choose the option that does not update unrelated packages.\\
\texttt{library(uafR)} fails right after install & RStudio has not refreshed the library. & Restart RStudio, then run \texttt{library(uafR)}.\\
No library detected & The chemical library argument was not supplied to \texttt{categorate()}. & Load \texttt{library\_data} and pass \texttt{chemical\_library = library\_data}.\\
Long database query & Live PubChem, KEGG, LOTUS, or other services are responding slowly. & Start with one or two compounds, keep \texttt{assay\_detail\_limit = 0}, use caching, and scale up after the smoke test works.\\
Missing output tables & A live lookup failed or returned incomplete source data. & Run \texttt{validateCategorateResult()} and inspect \texttt{ValidationSummary}, \texttt{TableQuality}, and \texttt{SourceDiagnostics}.\\
\bottomrule
\end{tabularx}
